\label{der:6-1}
实对称谱定理给$Q\T AQ=\diag(\lambda_i)$。对每个非零根作缩放$y_i=\frac{z_i}{\sqrt{|\lambda_i|}}$，零根方向保持不变，得到系数为$1,-1,0$的规范形；正交矩阵与这些非零缩放的乘积可逆，故为合同。

设一种规范形正系数个数为$p$，其正坐标子空间经可逆代换送回原向量空间，记为$U$；可逆代换保维数，故$\dim U=p$，且$f$在$U$上除零外严格正。若另一种规范形正系数为$p'$，将其非正坐标子空间也送回同一个原空间，记为$W$，则$\dim W=n-p'$，且$f$在$W$上不大于零。若$p>p'$，则$\dim U+\dim W>n$，由维数公式必有非零公共向量；在该向量上$f$既严格正又非正，矛盾。因此$p\le p'$，交换两种表示得$p'\le p$，故$p=p'$。对$-f$作相同论证得到负指数$q$也唯一。

可逆左右乘保持秩，规范形只有$p+q$个非零对角元，所以$\rank A=p+q$。若两实对称矩阵合同，惯性由上面论证相同；反之，它们分别合同于同一规范形$J$，设$P\T AP=J=Q\T BQ$，则$B=(PQ\inv)\T A(PQ\inv)$，因此合同。这里也验证了合同关系的反向构造。
